\documentclass[12pt,a4paper]{article}
\usepackage[utf8]{inputenc}
\usepackage[T1]{fontenc}
\usepackage[french]{babel}
\usepackage{geometry}
\geometry{margin=2.5cm}
\usepackage{graphicx}
\usepackage{booktabs}
\usepackage{hyperref}
\usepackage{amsmath}
\usepackage{enumitem}
\usepackage{xcolor}
\usepackage{float}
\usepackage{listings}
\usepackage{fancyhdr}
\usepackage{tcolorbox}

\pagestyle{fancy}
\fancyhead[L]{CANlock}
\fancyhead[R]{Session H26}
\fancyfoot[C]{\thepage}

\definecolor{navy}{HTML}{1B2A4A}
\hypersetup{colorlinks=true, linkcolor=navy, urlcolor=navy, citecolor=navy}

\title{\textcolor{navy}{\textbf{CANlock --- Système de Detection d'Intrusions sur Bus CAN}}\\[0.5em]
\large Guide complet du projet --- Session H26}
\author{Équipe CIA -- Universite Laval / Thales}
\date{Avril 2026}

\begin{document}
\maketitle

\begin{tcolorbox}[colback=navy!5, colframe=navy, title=\textbf{Résumé en 30 secondes}]
CANlock est un système de détection d'intrusions pour les véhicules connectés. Les véhicules modernes utilisent un réseau interne appelé \textbf{bus CAN} pour que leurs composants électroniques communiquent (freins, moteur, direction, etc.). Ce réseau n'a aucune sécurité native : un attaquant peut y injecter de faux messages pour prendre le contrôle du véhicule.

Notre solution : un modèle d'apprentissage profond (\textbf{CNN-LSTM Autoencoder}) qui apprend le comportement normal du réseau et détecte automatiquement les messages malveillants. Résultat : \textbf{99.5\% des attaques détectees} avec seulement \textbf{6.2\% de fausses alertes}.
\end{tcolorbox}

\tableofcontents
\newpage

% ============================================================================
\section{Contexte : pourquoi ce projet ?}
% ============================================================================

\subsection{Le bus CAN en bref}

Le \textbf{Controller Area Network (CAN)} est le protocole de communication standard dans l'industrie automobile depuis les années 1990. Dans un véhicule moderne, des dizaines d'unités électroniques (ECU) échangent en permanence des messages sur ce bus :

\begin{itemize}[noitemsep]
    \item L'ECU moteur envoie le régime et le couple
    \item L'ECU freinage envoie la pression sur les pédales
    \item Le tableau de bord lit ces messages pour afficher les informations
    \item Le système ABS réagit aux messages de freinage
\end{itemize}

\textbf{Le problème fondamental :} le protocole CAN a été conçu en 1986 pour la fiabilité, pas pour la sécurité. Il ne possède \textbf{aucun mécanisme d'authentification ni de chiffrement}. N'importe quel appareil connecté au bus peut envoyer n'importe quel message, et tous les ECU le recevront.

\subsection{La menace}

Un attaquant accédant au bus CAN (via le port OBD-II, un dongle Bluetooth compromis, ou le système d'infotainment) peut :
\begin{itemize}[noitemsep]
    \item \textbf{Injecter} de faux messages (ex: ``freiner a fond'')
    \item \textbf{Modifier} des messages en transit (ex: changer la vitesse affichee)
    \item \textbf{Supprimer} des messages (ex: empecher l'ABS de fonctionner)
    \item \textbf{Saturer} le bus pour bloquer toute communication
\end{itemize}

Ces attaques ont été demontrees dans la litterature sur des véhicules reels (Jeep Cherokee, Tesla, Toyota Prius).

\subsection{Objectif de CANlock}

\begin{tcolorbox}[colback=navy!5, colframe=navy]
Concevoir un \textbf{système de détection d'intrusions (IDS)} basé sur l'apprentissage profond, capable de détecter en temps réel les messages CAN malveillants avec :
\begin{itemize}[noitemsep]
    \item Un \textbf{recall} $\geq$ 99\% (ne rater presque aucune attaque)
    \item Un \textbf{taux de fausses alertes} faible (pour ne pas noyer l'operateur)
\end{itemize}
\end{tcolorbox}

% ============================================================================
\section{Les données}
% ============================================================================

\subsection{Source}

Les données proviennent d'un jeu de données de bus CAN collecte depuis un \textbf{poids lourd de marque Renault}. Le jeu de données contient environ \textbf{490 millions d'enregistrements} respectant le protocole \textbf{CAN J1939} (standard pour véhicules lourds).

\subsection{Structure d'un message CAN}

Chaque message CAN (128 bits) est compose de :
\begin{itemize}[noitemsep]
    \item \textbf{Identifier} (29 bits) : identifie le type de message (PGN) et l'ECU source
    \item \textbf{DLC} (4 bits) : longueur du champ de données en octets
    \item \textbf{Data field} (0 a 64 bits) : les données utiles (le ``payload'')
\end{itemize}

\textbf{Exemple concret :} un message avec l'identifiant 0x18FEF100 signifie ``PGN 65265 (Vehicle Speed) envoyé par l'ECU a l'adresse 0x00''. Le payload contient la vitesse encodée sur quelques bits.

\subsection{Le protocole J1939}

J1939 définit une hiérarchie de décodage :
\begin{itemize}[noitemsep]
    \item \textbf{PGN} (Parameter Group Number) : identifie le groupe de paramètres (ex: ``Electronic Engine Controller'')
    \item \textbf{SPN} (Suspect Parameter Number) : identifie un paramètre individuel dans le PGN (ex: ``Engine Speed'', ``Engine Torque'')
    \item Chaque SPN a une position en bits dans le payload, un facteur d'échelle et un offset pour convertir les bits bruts en valeur physique
\end{itemize}

\subsection{Stockage}

Les messages sont stockés dans une basé \textbf{PostgreSQL} (Docker, port 5435) avec le schéma :
\begin{verbatim}
Vehicle -> Session -> CanMessage (timestamp, can_identifier,
                                  length, payload)
PgnDefinition -> SpnDefinition (bit_start, bit_length,
                                 scale, offset)
\end{verbatim}

% ============================================================================
\section{Les attaques}
% ============================================================================

Pour tester notre IDS, nous generons des \textbf{attaques synthetiques réalistes} a partir de trafic CAN reel. Cinq types d'attaques couvrent les principales menaces documentees dans la litterature.

\subsection{DDoS (Denial of Service)}

\begin{tcolorbox}[colback=red!5, colframe=red!50!black]
\textbf{Objectif :} saturer le bus CAN pour empecher les communications normales.\\
\textbf{Comment :} envoi massif de messages sur un meme identifiant (3000 repetitions, intervalle de 0.0001s).\\
\textbf{Effet :} les messages légitimes sont retardes ou perdus, le véhicule peut devenir incontrolable.
\end{tcolorbox}

\subsection{Spoofing (Usurpation de payload)}

\begin{tcolorbox}[colback=red!5, colframe=red!50!black]
\textbf{Objectif :} injecter de fausses données en usurpant un ECU légitime.\\
\textbf{Comment :} pour 10\% des messages, modifie les octets du payload en appliquant un bruit gaussien.\\
\textbf{Effet :} le véhicule recoit des valeurs de capteur erronees (ex: fausse vitesse, faux couple moteur).
\end{tcolorbox}

\subsection{Replay (Rejeu)}

\begin{tcolorbox}[colback=red!5, colframe=red!50!black]
\textbf{Objectif :} capturer puis réémettre des messages légitimes a un autre moment.\\
\textbf{Comment :} 10\% des messages sont copiés et réinjectés avec 1 seconde de délai.\\
\textbf{Effet :} force une action passee (ex: rejeu d'une commande de deverrouillage).
\end{tcolorbox}

\subsection{Suspension (Bus-Off)}

\begin{tcolorbox}[colback=red!5, colframe=red!50!black]
\textbf{Objectif :} deconnecter un ECU du réseau.\\
\textbf{Comment :} génère des erreurs ciblées qui incrémentent le compteur d'erreurs de l'ECU victime (+8 par erreur). Au delà de 256, l'ECU se met en Bus-Off (mécanisme natif CAN).\\
\textbf{Effet :} l'ECU ciblé est isolé, ses messages disparaissent du bus.
\end{tcolorbox}

\subsection{Masquerade (Usurpation d'identité)}

\begin{tcolorbox}[colback=red!5, colframe=red!50!black]
\textbf{Objectif :} usurper l'identité d'un ECU pour prendre le contrôle de ses fonctions.\\
\textbf{Comment :} modifie l'adresse source (LSB) de l'identifiant CAN (8\% des messages).\\
\textbf{Effet :} l'attaquant parle ``au nom'' d'un ECU légitime (ex: se faire passer pour l'ECU moteur).
\end{tcolorbox}

\subsection{Paramètres d'attaque}

Les taux d'injection par défaut (2--3\%) étaient trop faibles : les séquences ``attaquées'' contenaient majoritairement du trafic normal, rendant l'apprentissage inefficace. Les paramètres optimaux :

\begin{table}[H]
\centering
\begin{tabular}{lccc}
\toprule
\textbf{Attaque} & \textbf{Paramètre} & \textbf{Defaut} & \textbf{Optimal} \\
\midrule
Spoofing & injection\_rate & 3\% & \textbf{10\%} \\
Masquerade & prob & 2\% & \textbf{8\%} \\
Replay & replay\_rate & 3\% & \textbf{10\%} \\
Suspension & suspend\_fraction & 2\% & \textbf{8\%} \\
DDoS & repetitions & 1000 & \textbf{3000} \\
\bottomrule
\end{tabular}
\caption{Paramètres d'attaque : valeurs par défaut vs optimales}
\end{table}

% ============================================================================
\section{Comment le modèle détecte les attaques}
% ============================================================================

\subsection{Le principe : détection par reconstruction}

Notre approche est \textbf{non supervisee} : le modèle n'a \textbf{jamais vu d'attaque} pendant l'entraînément. Il apprend uniquement a reconstruire du trafic CAN normal.

\begin{enumerate}
    \item \textbf{Entraînement :} on donne au modèle des milliers de séquences de trafic normal. Il apprend a les compresser puis les reconstruire fidèlement.
    \item \textbf{Detection :} face a un nouveau message, le modèle tente de le reconstruire. Si l'erreur de reconstruction est \textbf{élevée}, le message est ``hors distribution'' $\to$ probablement une \textbf{anomalie}.
\end{enumerate}

\textbf{Analogie :} comme un employe de banque qui voit des milliers de signatures normales chaque jour. Sans avoir jamais vu de faux, il détecte quand une signature ``ne ressemble pas'' a ce qu'il connait.

\subsection{Les 11 features}

Plutot que de décoder les messages CAN en valeurs physiques (SPN), ce qui est lent et perd de l'information, nous extrayons \textbf{11 features directement des messages bruts} :

\begin{table}[H]
\centering
\begin{tabular}{llp{5.5cm}}
\toprule
\textbf{Feature} & \textbf{Description} & \textbf{Pourquoi ?} \\
\midrule
\texttt{can\_id} & Identifiant CAN & Change lors d'une masquerade ou DDoS \\
\texttt{byte\_0..7} & 8 octets du payload & Changes lors d'un spoofing \\
\texttt{delta\_t} & Temps depuis le msg précédent & Anormal lors d'un replay (délai), suspension (gap) ou DDoS (trop rapide) \\
\texttt{freq} & Fréquence de cet ID dans les 50 derniers msgs & Explose lors d'un DDoS, chute lors d'une suspension \\
\bottomrule
\end{tabular}
\caption{Les 11 features et ce qu'elles détectent}
\end{table}

\textbf{Pourquoi pas le décodage SPN ?} L'approche initiale décodait chaque message en valeurs SPN via la basé J1939. Problèmes :
\begin{itemize}[noitemsep]
    \item Tres lent ($\sim$15 min pour 40k messages vs quelques secondes avec les payloads bruts)
    \item Le forward-fill et la normalisation lissaient les modifications des attaques
    \item Résultat : ROC-AUC de 0.52 (= hasard)
\end{itemize}

\subsection{Pre-traitement des données}

\begin{enumerate}[noitemsep]
    \item \textbf{Clipping des outliers} au 99e percentile sur \texttt{delta\_t} et \texttt{freq} (evite les spikes d'entraînément)
    \item \textbf{Normalisation} StandardScaler (fit uniquement sur les données normales)
    \item \textbf{Séquences temporelles} : fenêtre de 50 messages, stride de 25 (chevauchement)
\end{enumerate}

% ============================================================================
\section{Partitionnement des données}
% ============================================================================

Le partitionnement est \textbf{crucial} pour la fiabilité des résultats. Deux règles :

\begin{tcolorbox}[colback=navy!5, colframe=navy, title=\textbf{Règles de partitionnement}]
\begin{enumerate}
    \item \textbf{Split temporel} (pas de shuffle) : les données CAN sont des series temporelles. Melanger passe et futur causerait du \textit{data leakage}.
    \item \textbf{Attaques uniquement dans le test} : le modèle ne voit \textbf{jamais} d'attaque pendant l'entraînément ni la validation.
\end{enumerate}
\end{tcolorbox}

\textbf{Schema concret :}
\begin{itemize}[noitemsep]
    \item 500 000 messages charges depuis la BD, tries par timestamp
    \item 80\% premiers $\to$ \textbf{train + val} (uniquement normaux)
    \item 20\% derniers $\to$ \textbf{test} (on y applique les 5 attaques)
    \item Les séquences normales sont splittees : \textbf{75\% train / 15\% val / 10\% test}
    \item Le test final = 10\% normales + toutes les attaques $\to$ permet de mesurer FP et TP
\end{itemize}

Avec 500k messages, stride=25, seq\_len=50 : environ \textbf{12 000 séquences d'entraînément}, \textbf{2 400 de validation}, \textbf{5 768 de test} (1 601 normales + 4 167 attaquées).

% ============================================================================
\section{Les deux architectures évaluées}
% ============================================================================

Les deux modèles sont des \textbf{autoencodeurs} : ils compriment l'entrée en une représentation compacte (``bottleneck'') puis tentent de reconstruire l'entrée originale. La différence est dans comment ils compriment.

\subsection{RNN-VAE (Variational Auto-Encoder avec LSTM)}

\textbf{Idée :} apprendre une \textit{distribution} du trafic normal, pas juste une compression.

\textbf{Architecture :}
\begin{itemize}[noitemsep]
    \item \textbf{Encodeur :} LSTM bidirectionnel (64 hidden) $\to$ Dropout 0.2 $\to$ deux couches linéaires qui produisent $\mu$ et $\log\sigma^2$ (32 dimensions chacune)
    \item \textbf{Espace latent :} $z = \mu + \sigma \cdot \epsilon$ avec $\epsilon \sim \mathcal{N}(0, I)$ (reparameterization trick)
    \item \textbf{Décodeur :} linéaire 32$\to$128 $\to$ repeat sur la longueur de sequence $\to$ LSTM bidirectionnel $\to$ Dropout $\to$ linéaire $\to$ 11 features
\end{itemize}

\textbf{Loss :}
\begin{equation}
    \mathcal{L} = \underbrace{\text{MSE}(x, \hat{x})}_{\text{reconstruction}} + \underbrace{10^{-3} \cdot D_{KL}(\mathcal{N}(\mu, \sigma^2) \| \mathcal{N}(0, I))}_{\text{régularisation latente}}
\end{equation}

\textbf{152K paramètres} entrainables.

\subsection{CNN-LSTM Autoencoder (modèle retenu)}

\textbf{Idée :} utiliser des convolutions pour capter les patterns \textit{locaux} dans le payload, et un LSTM pour les dépendances \textit{temporelles}.

\textbf{Architecture :}
\begin{itemize}[noitemsep]
    \item \textbf{Encodeur CNN :} Conv1d(11$\to$32) + BN + ReLU + Pool(/2) $\to$ Conv1d(32$\to$64) + BN + ReLU + Pool(/2)
    \item \textbf{Encodeur LSTM :} LSTM bidirectionnel (32 hidden) + Dropout 0.2
    \item \textbf{Bottleneck :} linéaire 64$\to$32
    \item \textbf{Décodeur LSTM :} LSTM bidirectionnel (64 hidden)
    \item \textbf{Décodeur CNN :} ConvTranspose1d(128$\to$64) + BN + ReLU $\to$ ConvTranspose1d(64$\to$32) + BN + ReLU $\to$ Conv1d(32$\to$11) + interpolation
\end{itemize}

\textbf{Loss :}
\begin{equation}
    \mathcal{L} = \text{MSE}(x, \hat{x})
\end{equation}

\textbf{106K paramètres} entrainables.

\subsection{Pourquoi le CNN-LSTM est meilleur}

\begin{table}[H]
\centering
\begin{tabular}{lcc}
\toprule
\textbf{Aspect} & \textbf{RNN-VAE} & \textbf{CNN-LSTM} \\
\midrule
Patterns locaux (spoofing) & LSTM seul & \textbf{CNN} (convolutions) \\
Patterns temporels (DDoS, replay) & LSTM & \textbf{LSTM} \\
Stabilite entraînément & Trade-off MSE/KL & \textbf{MSE seul} \\
Fausses alertes (seuil recall 99.5\%) & 29.4\% & \textbf{6.2\%} \\
\bottomrule
\end{tabular}
\caption{Pourquoi le CNN-LSTM surpasse le RNN-VAE}
\end{table}

% ============================================================================
\section{Entraînement}
% ============================================================================

\subsection{Hyperparamètres}

Les deux modèles ont été entraînés avec les \textbf{memes hyperparamètres} pour une comparaison équitable :

\begin{table}[H]
\centering
\begin{tabular}{lc}
\toprule
\textbf{Hyperparamètre} & \textbf{Valeur} \\
\midrule
Messages charges & 500 000 \\
Batch size & 128 \\
Learning rate & $5 \times 10^{-4}$ \\
Epochs max & 200 \\
Sequence length & 50 \\
Stride & 25 \\
Early stopping patience & 20 epochs \\
Gradient clipping & 1.0 \\
Optimizer & Adam \\
Scheduler & ReduceLROnPlateau (factor=0.5, patience=5) \\
Seed & 42 \\
\bottomrule
\end{tabular}
\caption{Hyperparamètres d'entraînément (identiques pour les deux modèles)}
\end{table}

\subsection{Infrastructure et outils}

\begin{itemize}[noitemsep]
    \item \textbf{Matériel :} Intel Core 7 @ 1.8 GHz, CPU uniquement (pas de GPU)
    \item \textbf{Framework :} PyTorch Lightning
    \item \textbf{Suivi d'expériences :} MLflow (\texttt{http://localhost:5050})
    \item \textbf{Base de données :} PostgreSQL (Docker, port 5435)
    \item \textbf{Gestion des dépendances :} uv + pyproject.toml
    \item \textbf{Temps d'entraînément :} $\sim$30--45 min par modèle (CPU)
\end{itemize}

\subsection{Commande d'entraînément}

\begin{lstlisting}[basicstyle=\ttfamily\small, breaklines=true]
uv run train-models --model-type cnn_lstm \
    --limit 500000 --epochs 200 --batch-size 128 \
    --stride 25 --seq-len 50 --lr 5e-4 \
    --spoof-injection-rate 0.10 --masq-prob 0.08 \
    --replay-rate 0.10 --susp-suspend-fraction 0.08 \
    --ddos-repetitions 3000 --target-recall 0.995
\end{lstlisting}

% ============================================================================
\section{Detection d'anomalies : du modèle a la decision}
% ============================================================================

\subsection{Calcul de l'erreur de reconstruction}

Pour chaque sequence de test, on calcule l'erreur MSE entre l'entrée et la sortie du modèle :
\begin{equation}
    e_i = \frac{1}{T \times F} \sum_{t=1}^{T} \sum_{f=1}^{F} (x_{i,t,f} - \hat{x}_{i,t,f})^2
\end{equation}
ou $T=50$ (longueur de sequence) et $F=11$ (nombre de features).

\textbf{Intuition :} le modèle reconstruit bien ce qu'il connait (trafic normal $\to$ erreur faible). Face a une attaque, la reconstruction est mauvaise $\to$ erreur élevée.

\subsection{Choix du seuil}

L'erreur de reconstruction seule ne suffit pas : il faut un \textbf{seuil} au-dessus duquel on declare ``anomalie''. Deux stratégies :

\begin{enumerate}
    \item \textbf{Seuil F1-optimal :} maximise le F1-score (équilibre precision/recall). Bon pour un usage general.
    \item \textbf{Seuil recall-ciblé :} garantit un recall minimum (ex: 99.5\%). Indispensable pour la sécurité vehiculaire : mieux vaut quelques fausses alertes qu'une attaque ratée.
\end{enumerate}

\subsection{Pourquoi le recall est prioritaire}

\begin{itemize}[noitemsep]
    \item \textbf{Faux négatif} (attaque ratée) $\to$ conséquence potentiellement mortelle (prise de contrôle du véhicule)
    \item \textbf{Faux positif} (fausse alerte) $\to$ génère un log/alerte, mais pas de danger
    \item Sur un bus CAN avec des centaines de milliers de messages/minute, meme un FPR de 6\% génère beaucoup d'alertes $\to$ filtrage en aval nécessaire
\end{itemize}

% ============================================================================
\section{Résultats}
% ============================================================================

\subsection{Comparaison CNN-LSTM vs RNN-VAE}

Résultats avec le seuil recall $\geq$ 99.5\% :

\begin{table}[H]
\centering
\begin{tabular}{lcc}
\toprule
\textbf{Métrique} & \textbf{RNN-VAE} & \textbf{CNN-LSTM} \\
\midrule
ROC-AUC & 0.990 & \textbf{0.998} \\
FPR (fausses alertes) & 29.4\% & \textbf{6.2\%} \\
Recall (détection) & 99.5\% & 99.5\% \\
F1-Score & 0.940 & \textbf{0.988} \\
Accuracy & 91.0\% & \textbf{98.0\%} \\
\bottomrule
\end{tabular}
\caption{Comparaison des deux modèles (seuil recall $\geq$ 99.5\%)}
\end{table}

\textbf{Lecture :} les deux modèles détectent 99.5\% des attaques, mais le CNN-LSTM génère \textbf{5 fois moins de fausses alertes} (6.2\% vs 29.4\%).

\subsection{Matrices de confusion --- CNN-LSTM}

\begin{table}[H]
\centering
\begin{tabular}{l|cc|cc}
\toprule
& \multicolumn{2}{c|}{\textbf{Seuil F1-optimal}} & \multicolumn{2}{c}{\textbf{Seuil Recall $\geq$ 99.5\%}} \\
& Pred. Normal & Pred. Anomalie & Pred. Normal & Pred. Anomalie \\
\midrule
Reel Normal (1601) & 1559 & 42 & 1502 & 99 \\
Reel Anomalie (4167) & 45 & 4122 & 21 & 4146 \\
\midrule
Recall & \multicolumn{2}{c|}{98.9\%} & \multicolumn{2}{c}{\textbf{99.5\%}} \\
FPR & \multicolumn{2}{c|}{\textbf{2.6\%}} & \multicolumn{2}{c}{6.2\%} \\
\bottomrule
\end{tabular}
\caption{Matrices de confusion du CNN-LSTM}
\end{table}

\textbf{Lecture :} avec le seuil recall $\geq$ 99.5\%, le modèle ne rate que \textbf{21 attaques sur 4167}. Il génère 99 fausses alertes sur 1601 séquences normales.

\subsection{Matrices de confusion --- RNN-VAE}

\begin{table}[H]
\centering
\begin{tabular}{l|cc|cc}
\toprule
& \multicolumn{2}{c|}{\textbf{Seuil F1-optimal}} & \multicolumn{2}{c}{\textbf{Seuil Recall $\geq$ 99.5\%}} \\
& Pred. Normal & Pred. Anomalie & Pred. Normal & Pred. Anomalie \\
\midrule
Reel Normal (1601) & 1486 & 115 & 1131 & 470 \\
Reel Anomalie (4167) & 134 & 4033 & 21 & 4146 \\
\midrule
Recall & \multicolumn{2}{c|}{96.8\%} & \multicolumn{2}{c}{99.5\%} \\
FPR & \multicolumn{2}{c|}{7.2\%} & \multicolumn{2}{c}{29.4\%} \\
\bottomrule
\end{tabular}
\caption{Matrices de confusion du RNN-VAE}
\end{table}

\textbf{Lecture :} le RNN-VAE atteint le meme recall, mais génère \textbf{470 fausses alertes} (vs 99 pour le CNN-LSTM).

\subsection{Evolution des performances : comment on y est arrive}

L'optimisation a été itérative. Chaque étape a apporte une amélioration mesurable :

\begin{table}[H]
\centering
\small
\begin{tabular}{lcccp{5cm}}
\toprule
\textbf{Étape} & \textbf{ROC-AUC} & \textbf{FPR} & \textbf{Recall} & \textbf{Ce qui a change} \\
\midrule
Départ (SPN, 50k) & 0.52 & 1.00 & 0.997 & Décodage SPN, tout classe anomalie \\
Payloads bruts & 0.54 & 1.00 & 0.997 & 8 octets seulement, pas assez \\
+ CAN ID + temps & 0.67 & 0.99 & 0.997 & 11 features, debut de discrimination \\
+ 200k messages & 0.70 & 0.68 & 0.917 & Plus de données d'entraînément \\
+ 500k, stride=10 & 0.92 & 0.29 & 0.917 & Beaucoup plus de séquences \\
+ Taux attaque x3 & 0.996 & 0.037 & 0.979 & Attaques plus visibles \\
\textbf{Config finale} & \textbf{0.998} & \textbf{0.062} & \textbf{0.995} & Epochs 200, patience 20 \\
\bottomrule
\end{tabular}
\caption{Evolution des performances du CNN-LSTM (chaque ligne = une experience)}
\end{table}

% ============================================================================
\section{Architecture logicielle}
% ============================================================================

\subsection{Structure du code}

\begin{verbatim}
src/canlock/
    attacks/              # Generation d'attaques
        attack_base.py    #   Classe abstraite + utilitaires bit-level
        attack_ddos.py    #   Attaque DDoS
        spoofing_attack.py#   Spoofing (bruit gaussien)
        replay_attack.py  #   Replay (rejeu avec délai)
        suspension_attack.py # Suspension (Bus-Off)
        masquerade_attack.py # Masquerade (usurpation ID)
    data/
        can_data_module.py #  DataModule PyTorch Lightning
    db/
        models.py         #  ORM SQLModel (Vehicle, Session,
                          #    CanMessage, PGN, SPN)
        database.py       #  Connexion PostgreSQL
    models/
        cnn_lstm_autoencoder.py  # CNN-LSTM (modèle retenu)
        rnn_vae.py               # RNN-VAE
        anomaly_detector.py      # Évaluation (seuil, métriques)
    training/
        train_models.py   #  Pipeline d'entraînément complet (CLI)
    décoder.py            #  Décodeur SPN/PGN J1939
    cli.py                #  Points d'entrée CLI

scripts/
    attack_synthesis.py   # Generation de datasets attaques
    attack_inspect.py     # Inspection/validation
    attack_visualize.py   # Visualisation
    generate_poster_figures.py     # Figures pour l'affiche
    generate_architecture_figures.py # Diagrammes d'architecture

notebooks/
    rnn_vae_model.ipynb    # Expérimentation RNN-VAE
    cnn_lstm_model.ipynb   # Expérimentation CNN-LSTM
    can_data_analysis.ipynb # Analyse exploratoire
\end{verbatim}

\subsection{Services Docker}

\begin{lstlisting}[basicstyle=\ttfamily\small, breaklines=true]
# docker-compose.yml
services:
  postgres-backend:  # BDD CAN (port 5435)
  mlflow:            # Suivi d'expériences (port 5050)
\end{lstlisting}

\subsection{CLI : commandes disponibles}

\begin{lstlisting}[basicstyle=\ttfamily\small, breaklines=true]
uv run train-models     # Entraînement des modèles
uv run download-heavy-truck-data  # Telechargement données
uv run analyze-session  # Analyse d'une session CAN
\end{lstlisting}

% ============================================================================
\section{Conclusion}
% ============================================================================

\begin{tcolorbox}[colback=navy!5, colframe=navy, title=\textbf{Résultats cles}]
\begin{itemize}[noitemsep]
    \item Le \textbf{CNN-LSTM Autoencoder} détecte \textbf{99.5\% des intrusions} avec \textbf{6.2\% de fausses alertes}
    \item ROC-AUC de \textbf{0.998} (quasi-parfait)
    \item Detecte les \textbf{5 types d'attaques} : DDoS, spoofing, masquerade, replay, suspension
    \item Entraînement \textbf{non supervise} (pas besoin d'exemples d'attaques pour entraînér)
    \item \textbf{11 features} extraites des messages bruts (pas de décodage SPN couteux)
    \item Surpasse le RNN-VAE : \textbf{5x moins de fausses alertes} pour le meme recall
\end{itemize}
\end{tcolorbox}

\subsection{Travaux futurs}

\begin{itemize}[noitemsep]
    \item Validation sur des \textbf{attaques réelles} (pas uniquement synthetiques)
    \item Optimisation pour \textbf{inférence en temps reel} sur hardware embarqué (microcontrôleur)
    \item Ajout d'attaques plus sophistiquées (fuzzing, attaques graduelles, attaques zero-day)
    \item Évaluation sur d'autres jeux de données CAN (véhicules légers, protocole CAN-FD)
    \item Reduction du FPR par filtrage contextuel post-détection
\end{itemize}

% ============================================================================
\section*{Références}
% ============================================================================

\begin{enumerate}[noitemsep, label={[\arabic*]}]
    \item Rajapaksha, S. et al. (2023). Beyond vanilla: Improved autoencoder-based ensemble in-vehicle IDS. \textit{J. Information Security and Applications}, 77.
    \item Balaji, P. et al. (2022). CANLite: Anomaly Detection in CAN with Multitask Learning. \textit{IEEE VTC2022-Spring}.
    \item Alkhatib, N. et al. (2022). CAN-BERT do it? CAN IDS based on BERT. \textit{IEEE AICCSA}.
    \item Gao, F. et al. (2024). Multi-Attack IDS for In-Vehicle CAN-FD Messages. \textit{Sensors}, 24(11).
    \item Hoang, T.-N. \& Kim, D. (2022). Semi-supervised convolutional adversarial autoencoders for in-vehicle intrusion détection. \textit{Vehicular Communications}, 38.
    \item Mansourian, P. et al. (2023). Deep Learning-Based Anomaly Detection for CAVs Using Spatiotemporal Information. \textit{IEEE TITS}, 24(12).
    \item Xiao, J. et al. (2023). Robust anomaly-based IDS for in-vehicle network by graph neural network framework. \textit{Applied Intelligence}, 53(3).
    \item Cao, J. et al. (2024). Anomaly Detection for In-Vehicle Network Using Self-Supervised Learning. \textit{IEEE TITS}, 25(7).
\end{enumerate}

\end{document}
